\subsection{赋范空间的情形}

设 \(\mathbf{E}\) 是一个赋范空间。对偶 \(\mathbf{E}'\)（\(\mathbf{E}\) 的）上的强拓扑由范数

\[
\left\| x ^ {\prime} \right\| = \sup _ {x \in \mathrm{E}, \| x \| \leqslant 1} \left| \langle x, x ^ {\prime} \rangle \right|,\tag{1}
\]

定义，而 E 的强对偶是一个 Banach 空间（III，第 24 页，推论 2）。于是 E 的双对偶 E'' 也是一个 Banach 空间，其范数由

\[
\left\| x ^ {\prime \prime} \right\| = \sup _ {x ^ {\prime} \in \mathrm{E} ^ {\prime}, \left\| x ^ {\prime} \right\| \leqslant 1} \left| \left\langle x ^ {\prime}, x ^ {\prime \prime} \right\rangle \right|.\tag{2}
\]

由 IV 第 7 页的命题 8 (i)，典范线性映射 \(c_{\mathrm{E}}: \mathrm{E} \to \mathrm{E}''\) 是一个等距。此后，我们将把 E 等同于它的双对偶 \(\mathrm{E}''\) 的一个赋范子空间。

命题 5. --- 设 E 是一个赋范空间，E' 是它的对偶，E'' 是它的双对偶。E'' 中的（闭）单位球是 E 中单位球 B 对于弱拓扑 \(\sigma(\mathrm{E}'', E')\) 的闭包。由公式 (1) 与 (2)，E'' 中的单位球是 B 的双极集 B°°。于是命题 5 由双极定理（II，第 45 页，推论 3）得出。

注. --- Banach 空间 E 在它的双对偶 E'' 中对于强拓扑是闭的，但对于弱拓扑是稠密的（命题 5）。

为使一个赋范空间是自反的，必要且充分的是它是半自反的；因为 E 的初始拓扑总是由 \(E''\) 的强拓扑诱导的。于是（IV，第 15 页的）定理 1 蕴含以下结果：

命题 6. --- 为使赋范空间 E 是自反的，必要且充分的是 E 中的单位球对于弱化拓扑 \(\sigma(E, E')\) 是紧的。

我们注意到，自反赋范空间是完备的，因而是 Banach 空间，且由 IV 第 16 页的命题 4，它的对偶是自反 Banach 空间。

命题 7. --- 设 E 是一个自反 Banach 空间，M 是 E 的一个闭向量子空间。则 M 与 E/M 都是自反 Banach 空间。

设 \(E'\) 是 E 的对偶，\(M^{\circ}\) 是 M 在 \(E'\) 中的正交集。作为赋范空间，我们可以把空间 \(E'/M^{\circ}\) 等同于 M 的对偶 \(M'\)（IV，第 9 页，命题 10）。由于 M 是半自反的（IV，第 16 页，推论），它是自反的，从而 \(E'/M^{\circ}\) 也是自反的；类似地，\(M^{\circ}\) 是自反的，它的对偶 \(E/M^{\circ\circ} = E/M\) 亦然。

例. --- 1) 设 \(\ell^{\infty}(\mathbf{N})\) 表示由有界序列 \(\mathbf{x} = (x_n)_{n\in \mathbf{N}}\)

（的数列）组成的 Banach 空间，其范数为

\[
\| \mathbf {x} \| = \sup _ {n \in \mathbf {N}} | x _ {n} | \quad (\mathrm{I}, \text {p.} 4).\tag{3}
\]

设 \(c_{0}(\mathbf{N})\) 是 \(\ell^{\infty}(\mathbf{N})\) 中由趋于 0 的序列组成的闭向量子空间。最后，设 \(\ell^{1}(\mathbf{N})\) 是可求和序列组成的向量空间，赋以范数

\[
\left\| \mathbf {x} \right\| _ {1} = \sum_ {n \in \mathbf {N}} \left| x _ {n} \right|.\tag{4}
\]

我们可以证明（IV，第 47 页，习题 1），\(c_0(\mathbf{N})\) 的对偶可以等同于 \(\ell^1 (\mathbf{N})\)，使得我们有

\[
\langle \mathbf {x}, \mathbf {x} ^ {\prime} \rangle = \sum_ {n \in \mathbf {N}} x _ {n} x _ {n} ^ {\prime}\tag{5}
\]

（对所有 \(\mathbf{x} \in c_0(\mathbf{N})\) 与 \(\mathbf{x}' \in \ell^1(\mathbf{N})\)）。类似地，\(\ell^1(\mathbf{N})\) 的对偶可以等同于 \(\ell^\infty(\mathbf{N})\)，使得关系式 (5) 对所有 \(\mathbf{x} \in \ell^1(\mathbf{N})\) 与所有 \(\mathbf{x}' \in \ell^\infty(\mathbf{N})\) 成立。因此 \(\ell^\infty(\mathbf{N})\) 是 \(c_0(\mathbf{N})\) 的双对偶，且后一空间不是自反的。

* 2) 每个 hilbertian 空间都是自反 Banach 空间（V，第 17 页）。*

* 3) 设 X 是一个 Hausdorff 拓扑空间，\(\mu\) 是 X 上的一个复测度。对每个实数 \(p > 1\)，Banach 空间 \(\mathrm{L}^p (\mathrm{X},\mu)\) 是自反的，且它的对偶可以等同于 \(\mathrm{L}^q (\mathrm{X},\mu)\)，其中 \(p^{-1} + q^{-1} = 1\)（INT，V，第二版，§ 5，第 8 目与 IX，§ 1，第 10 目）。*

